\section{Background}
\label{sec:background-ra}

Compensating transactions and the saga pattern give the long-running
analog to atomic transactions for workflows whose commit boundary is
weaker than serializable schedules~\cite{garciamolina1987sagas}. A saga is a
sequence of forward steps each paired with a compensating step that
semantically undoes the forward step. Workflow nets extend the model
with timed arcs and soundness criteria that decide whether every
reachable marking has a path to a final marking. The category of effects
reversible by inverse executors (file rename, signal stop / continue,
network policy reload) is a strict subclass of the compensating-action
literature; the reversible-action substrate inherits the saga
vocabulary without the workflow-net soundness obligations, which would
require a timed transition system the substrate does not provide.

Capability revocation in operating systems has the closest structural
analog. EROS, seL4, and Capsicum each treat revocation as a typed
operation on a capability token: the capability becomes unusable after
revocation, and the revocation event is itself a checkable artifact.
Bounded executive action shares the revocation semantics and adds a
constitutional admission layer above it: the act of revoking is not
just an OS operation but a receipt whose admission predicates include
the bilateral-cosignature requirement when the action class is
destructive. The OS-capability literature solved revocation at the
machine boundary; the substrate lifts the discipline to the cross-
organizational boundary where the receipt graph crosses kernels.

Endpoint detection and response engines have shipped TTL fields and
manual rollback workflows for a decade. Vendor schemas record an action
identifier, a target host, an action class, a TTL duration, and a
rollback handle; auditors inspect the rollback handle by hand when
something goes wrong. No deployed engine reduces TTL enforcement to a
typed invariant, and no engine documents a rollback receipt admission
predicate distinct from the original execution receipt's admission
predicate. The contribution is not a new engine; it is a typed
formulation of the four-receipt grammar (request, execution, rollback,
acknowledgement) and an admission predicate stack that the existing
engines could in principle adopt.

The prior substrate \cite{programmableSovereigntyAnonymous} supplies
four primitives the reversible-action discipline depends on. The treaty
intersection theorem
\thm{treaty_admission_iff_predicate_intersection} gives the bilateral
reduction. The ladder-floor stability theorem
\thm{treaty_admission_stable_under_ladder_floor} gives the receipt-
backed floor under which destructive actions are inadmissible. The
amendment refinement bridge
\thm{amendment_admissible_iff_backward_refinement} gives the type-level
invariant on enactment. The essential-predicate chain theorem
\thm{essential_preserved_chain} gives the trajectory-invariant
composition over amendment sequences. The composition theorem added
here is the response-side analog of the fourth primitive: TTL-bounded
amendments compose with the chain machinery the prior substrate already
established.

The Anthropic alignment-evaluation literature on alignment-faking and
situational-awareness is a non-substrate motivation. A scheming agent
that behaves under evaluation and differently in deployment leaves no
checkable artifact unless the runtime emits one. A typed reversible-
action discipline binds the agent's destructive actions to a bilateral
admission predicate the deploying organization participates in, so the
artifact cannot be retroactively denied. The substrate does not solve
scheming; it makes the consequence of an agent's destructive action
recordable in the same canonical form the prior substrate's receipts
\cite{programmableSovereigntyAnonymous} already take.
